\chapter{Introduction}
\ChapterLabel{Introduction}
\index{introduction}
\index{Chapter!Introduction}
\index{normative language}
\index{priority}
\index{source of truth}
\index{V-model}

\section{How to Use This Chapter}
State the document purpose, scope, audience, conventions, and relationship to
other project artifacts. Replace bracketed text with project-specific content.
Keep binding obligations in the Requirements chapter.

\section{Purpose}
This document is the authoritative engineering specification for
\ProjectName. It records why the system is needed, what it shall do, how it is
structured, how it will be tested, and what evidence is required before
release. It is readable by people and structured for use by engineering and
\gls{ai}-assisted development tools.

\begin{dsnNote}[Central lesson]
A vague request produces unpredictable software. A precise, traceable
specification allows engineers to generate, review, test, and, when necessary,
regenerate a conforming implementation.
\end{dsnNote}

\section{Scope}
Define the system boundary, included capabilities, exclusions, operating
environment, major constraints, and the release covered by this document.

\ProjectName\ is a system that controls the opening and closing of a
curtain. The current decision is that a user commands the curtain through
both a physical controller and a phone application. Stakeholders are identified in
\hyperref[Section:Stakeholders]{the System Context}. Included capabilities,
exclusions, target environment, external dependencies, and the release covered by this document
are TBD and require supporting project evidence. Open context items are
listed in \ChapterRef{SystemContext}.

\section{Audience}
Identify the people who create, approve, implement, test, operate, maintain,
audit, or use the system. State which sections are most relevant to each role.

\textit{Generic example: The primary audiences are customer representatives,
systems engineers, architects, developers, testers, operators, maintainers,
security reviewers, and project approvers. One person may perform several
roles on a small project.}

\section{Normative Language and Priority}
Use \textbf{shall} for every binding requirement. Do not use \textbf{should} as
guidance in normative statements. Assign each requirement one priority:
\textbf{\gls{must}}, \textbf{\gls{should}}, \textbf{\gls{could}}, or
\textbf{\gls{would}}. Priority expresses delivery importance; it does not change
whether a statement containing \textbf{shall} is binding for its assigned
implementation version.

\section{Document as Source of Truth}
Requirements alone do not define a reproducible implementation. The complete
source of truth consists of stakeholder needs, user stories, requirements, use
cases, architecture and interface decisions, implementation constraints,
tests, verification evidence, validation evidence, and release criteria.
Generated code is an implementation artifact governed by this specification.

\section{Initializing a New Project from This Template}
\label{Section:InitializeNewProject}
\index{template!project initialization}
\index{AI agent!project initialization}
Create a project-specific architecture book from a copy of this template; do
not place project-specific content in the master template. Before creating or
renaming project-owned files, confirm the project name, system boundary,
repository location, document owner, and project file namespace. Preserve
tool-mandated discovery filenames, such as \texttt{AGENTS.md}, when an AI tool
requires them.

Gather the approved project evidence and identify missing scope, stakeholders,
interfaces, constraints, quality attributes, risks, verification evidence,
validation evidence, and approvals before replacing placeholders. Replace only
content supported by that evidence. Update the document metadata, title
material, contact information, and document history; then develop stakeholder
needs, user stories, requirements, use cases, architecture, tests,
traceability, verification, validation, and release content in dependency
order. Preserve stable identifiers and reciprocal mappings throughout the
document.

Human authors and AI agents shall follow the detailed Initialize a Project
workflow, repository-bootstrap procedure, naming rules, and completion
checklist in \ChapterRef{AIAgentGuide}. After initialization, compile the
complete document, review the build log, and visually inspect the generated
PDF. The master template remains the reusable reference for later projects.

To initialize a project with an AI agent, replace the bracketed project name
and give the agent the following request:

\begin{verbatim}
Read Section 1.7, "Initializing a New Project from This Template,"
and the authoritative AI Agent Usage Guide in tgmAIAgentGuide.tex.
Initialize this copy of the architecture template for [Project Name].
Before editing, ask me to confirm the project file namespace, which
AI tools this repository will use, and which domain tailoring profiles
apply. Review the repository and supplied project evidence, then
produce a gap list of missing scope,
stakeholders, interfaces, constraints, quality attributes, risks,
evidence, and approvals. Do not invent project facts or replace
unsupported placeholders. After I review the gap list, follow the
documented initialization workflow and dependency order. Create only
the applicable AI bootstrap files. Create SKILL.md only if I confirm
that the repository is intended to provide a reusable AI skill. Update
document history, compile the complete document, inspect the build
log, and visually inspect the resulting PDF. Report unresolved
decisions that require human approval.
\end{verbatim}

\section{Template Tailoring and Domain Profiles}
\label{Section:DomainProfiles}
\index{domain tailoring profile}
\index{template!tailoring}
The core chapters form a generic software-architecture template. A domain
tailoring profile augments those chapters with concerns, artifacts, diagrams,
and review criteria for a particular class of system. It does not replace a
core chapter or create a second authoritative location for project facts.
Requirements, interfaces, architecture decisions, tests, and evidence remain
in their normal core chapters and link to one another through their stable IDs.

Select applicable profiles during project initialization. More than one
profile may apply. Record each applicability decision and rationale in
\ChapterRef{DevelopmentPlan}, retain applicable profile appendices as
authoritative tailoring guidance, and exclude an inapplicable profile from an
instantiated project only after recording the decision. If profile guidance
conflicts, identify the conflict and obtain an authorized human decision before
editing controlled project content.

\small
\begin{tabularx}{\linewidth}{L{4.0cm}L{3.8cm}Y}
\caption{Domain Profile Applicability\label{Table:DomainProfileApplicability}}\\
\toprule
\textbf{Profile} & \textbf{Applicability} & \textbf{Decision and Rationale}\\
\midrule
\endfirsthead
\multicolumn{3}{l}{\small\itshape Domain Profile Applicability (continued)}\\
\toprule
\textbf{Profile} & \textbf{Applicability} & \textbf{Decision and Rationale}\\
\midrule
\endhead
\bottomrule
\endlastfoot
\hyperref[Chapter:EmbeddedSystemsProfile]{Embedded Systems} &
Applicable &
The system controls the physical movement of a curtain. Decision record:
\hyperref[Table:ProjectDomainProfiles]{project profile decisions}.\\
\textit{[Additional profile]} &
\textit{[Applicable / Not Applicable / TBD]} &
\textit{[Reason and relationship to other profiles]}\\
\end{tabularx}
\normalsize

When the Embedded Systems profile is Applicable or TBD, read
\ChapterRef{EmbeddedSystemsProfile} during initialization and for every task
affected by hardware, firmware, real-time behavior, constrained resources,
physical interfaces, safety, security, manufacturing, or field service.

\section{Using This Template with AI Agents}
\index{AI agent!architecture workflow}
AI agents such as OpenAI Codex and Anthropic Claude Code can use the LaTeX
sources as structured project context. The agent shall begin with
\ChapterRef{AIAgentGuide}, select the applicable task-routing entry, and read
only the authoritative files identified for that task. It shall expand the
context only when a referenced dependency or unresolved inconsistency requires
it. This keeps routine work focused without weakening traceability.

The AI Agent Usage Guide contains the task-routing table, context packs,
mandatory rules, standard workflows, completion checklists, human-review
boundaries, and reusable prompts. Record material AI-assisted decisions in
\ChapterRef{AIPromptLog}.

The following minimal request starts an AI-assisted task:

\begin{verbatim}
Read tgmAIAgentGuide.tex first. Identify the applicable task and
load only the files listed for it. Follow the template's stable-ID,
priority, status, traceability, history, build, and review rules.
Use only supplied project evidence. Report missing information and
decisions that require human approval. Complete the requested task,
compile the document, and visually inspect the affected PDF pages.
\end{verbatim}

\section{V-Model Organization}
Figure~\ref{Figure:VModel} summarizes the \gls{vmodel} organization used by
this template. The left side moves from stakeholder needs toward detailed
design. The bottom represents implementation. The right side integrates the
system and gathers evidence, pairing acceptance with stakeholder needs, system
verification with requirements, integration tests with architecture, and unit
tests with detailed design.

This organization is compatible with the system and software life-cycle
process frameworks in \cite{ISO15288,ISO12207}; those standards do not require
a particular life-cycle model.

\begin{figure}[htbp]
\centering
\begin{tikzpicture}[
  stage/.style={
    draw=dsnBlue,
    fill=dsnLightBlue,
    rounded corners=2pt,
    minimum width=2.7cm,
    minimum height=0.82cm,
    align=center,
    font=\small
  },
  test stage/.style={
    stage,
    draw=dsnGreen,
    fill=dsnGreen!10
  },
  implementation/.style={
    stage,
    draw=dsnGray,
    fill=dsnGray!10
  },
  lifecycle/.style={
    -{Latex[length=2.2mm]},
    line width=0.7pt
  },
  relationship/.style={
    draw=dsnGray,
    dashed,
    line width=0.6pt
  },
  relationship label/.style={
    fill=white,
    inner sep=1.5pt,
    font=\scriptsize\itshape,
    text=dsnGray
  }
]
  \node[stage] (needs) at (-5.6,5.0) {Stakeholder Needs};
  \node[stage, minimum width=3.2cm] (requirements) at (-4.35,3.75) {System and Software\\Requirements};
  \node[stage] (architecture) at (-3.1,2.5) {Architecture};
  \node[stage] (design) at (-2.2,1.25) {Detailed Design};
  \node[implementation] (implementation) at (0,0) {Implementation};

  \node[test stage] (unit) at (2.2,1.25) {Unit Tests};
  \node[test stage] (integration) at (3.1,2.5) {Integration Tests};
  \node[test stage, minimum width=3.0cm] (verification) at (4.35,3.75) {System Verification};
  \node[test stage, minimum width=2.9cm] (validation) at (5.6,5.0) {Stakeholder Validation};

  \draw[lifecycle, draw=dsnBlue]
    (needs) -- (requirements) -- (architecture) -- (design) -- (implementation);
  \draw[lifecycle, draw=dsnGreen]
    (implementation) -- (unit) -- (integration) -- (verification) -- (validation);

  \draw[relationship] (needs) -- node[relationship label, above=1pt] {acceptance} (validation);
  \draw[relationship] (requirements) -- node[relationship label, above=1pt] {compliance} (verification);
  \draw[relationship] (architecture) -- node[relationship label, above=1pt] {integration} (integration);
  \draw[relationship] (design) -- node[relationship label, above=1pt] {correctness} (unit);

\end{tikzpicture}

\caption{V-model relationships used by this specification.}
\label{Figure:VModel}
\end{figure}
